\documentclass[a4paper,11pt]{article}
\usepackage[T2A]{fontenc}
\usepackage[utf8]{inputenc}
\usepackage[russian]{babel}
\usepackage{amsmath,amssymb,mathtools}
\renewcommand{\sfdefault}{cmss}
\renewcommand{\familydefault}{\sfdefault}
\usepackage{geometry}
\geometry{left=21mm,right=21mm,top=17mm,bottom=18mm}
\usepackage{microtype}
\usepackage{tikz}
\usepackage{tabularx,booktabs,array}
\usepackage{hyperref}
\usepackage{enumitem}
\usepackage{parskip}
\usepackage{fancyhdr}
\usepackage{xcolor}
\definecolor{accent}{HTML}{145E65}
\definecolor{ink}{HTML}{24282A}
\definecolor{muted}{HTML}{667277}
\definecolor{wrong}{HTML}{A32B32}
\hypersetup{colorlinks=true,linkcolor=accent,urlcolor=accent,pdfauthor={Лёвин Артём Александрович},pdftitle={Ревью домашней работы — Марк Гукин — 29.10.26}}
\setlength{\parindent}{0pt}
\setlength{\parskip}{3pt}
\setlength{\abovedisplayskip}{6pt}
\setlength{\belowdisplayskip}{6pt}
\setlength{\abovedisplayshortskip}{3pt}
\setlength{\belowdisplayshortskip}{3pt}
\linespread{1.04}
\setlist[itemize]{leftmargin=*,topsep=2pt,itemsep=2pt}
\pagestyle{fancy}
\fancyhf{}
\renewcommand{\headrulewidth}{0pt}
\fancyfoot[C]{\footnotesize\color{muted}Лёвин Артём Александрович эксклюзивно для Марка Гукина}
\fancyfoot[R]{\footnotesize\color{muted}\thepage}
\newcommand{\minor}[1]{{\color{muted}\small #1}}
\newcommand{\err}[1]{\textcolor{wrong}{\textbf{Ошибка:}} #1}
\newcommand{\keyidea}[1]{\textcolor{accent}{\textbf{Ключевая идея:}} #1}
\newcommand{\result}[1]{\textcolor{accent}{\textbf{Итог:}} #1}
\newcommand{\subfield}[1]{\vspace{1pt}{\color{accent}\bfseries #1}\par\vspace{-2pt}}
\newcommand{\taskhead}[2]{%
  \phantomsection\label{task:#1}%
  {\small\color{muted}БОЛЬШОЕ ПОВТОРЕНИЕ / ЗАДАНИЕ #1}\par
  \vspace{5pt}{\LARGE\bfseries\color{ink}#2}\par
  \vspace{2pt}{\color{accent}\rule{\linewidth}{0.6pt}}\par\vspace{5pt}%
}
\begin{document}

\thispagestyle{empty}
\vspace*{29mm}
\begin{center}
{\small\color{accent}\bfseries МАТЕМАТИКА / ПРОФИЛЬНЫЙ ЕГЭ}\par
\vspace{18mm}
{\Huge\bfseries Ревью домашней работы}\par
\vspace{17mm}
{\Large Марк Гукин}\par
\vspace{22mm}
{\large 29 октября 2026 года}\par
\vspace{15mm}
{\normalsize Лёвин Артём Александрович}\par
\end{center}
\vfill
\begin{center}
\begin{tikzpicture}
\draw[accent!45,line width=.85pt] (-4,0) .. controls (-2.0,.7) and (-1.2,-.55) .. (0,.05) .. controls (1.6,.8) and (2.4,-.3) .. (4,.2);
\end{tikzpicture}
\end{center}
\clearpage

{\small\color{muted}29.10.2026 / БОЛЬШОЕ ПОВТОРЕНИЕ}\par
\vspace{5pt}
{\LARGE\bfseries\color{ink}Сводная таблица}\par
\vspace{7pt}
В таблице перечислены задания, требующие исправления или завершения. Номер задания ведёт к подробному разбору.
\vspace{9pt}

\begingroup
\small
\renewcommand{\arraystretch}{1.36}
\begin{tabularx}{\linewidth}{@{}>{\raggedright\arraybackslash}p{9mm}>{\raggedright\arraybackslash}p{22mm}>{\raggedright\arraybackslash}X>{\raggedright\arraybackslash}p{31mm}>{\raggedright\arraybackslash}p{34mm}@{}}
\toprule
\textbf{№} & \textbf{Тема} & \textbf{Суть ошибки} & \textbf{Твой ответ} & \textbf{Правильный ответ}\\
\midrule
\hyperref[task:9]{9} & Дробная степень & Потерян отрицательный корень & $8$ & $-8;\ 8$\\
\hyperref[task:12]{12} & Логарифмы & Неверное устранение деления на 3 & $6$ & $\frac85$\\
\hyperref[task:17]{17} & Степенная зависимость & Решение не доведено до отношения температур & Не указан & $\frac32$ раза\\
\hyperref[task:18]{18} & Логарифмы & В формуле корней вместо $7$ записано $4$ & $\frac{4+\sqrt{41}}2$ & $\frac{7+\sqrt{41}}2$\\
\hyperref[task:20]{20} & Двойной угол & Не учтён знак синуса в IV четверти & \shortstack[l]{$\sin2t$\\$=\frac{24}{25}$} & \shortstack[l]{$\sin2t$\\$=-\frac{24}{25}$}\\
\bottomrule
\end{tabularx}
\endgroup
\vspace{11pt}
\minor{В задании №20 второе найденное значение, $\cos 2t=-\frac7{25}$, верно.}
\clearpage

\thispagestyle{fancy}
\vspace*{\fill}
\begin{center}
{\Huge\color{accent}\bfseries Разбор ошибок}
\end{center}
\vspace*{\fill}
\clearpage

\taskhead{9}{Дробная степень: два корня}
\subfield{Текст задания}
Реши в действительных числах уравнение
\[
 x^{\frac43}=16.
\]
\subfield{Твоё решение и ответ}
В записи использован переход от $x^{4/3}=2^4$ к возведению обеих частей в степень $3/4$, после чего получено $x=8$. \textbf{Твой ответ:} $8$.
\subfield{Где именно возникает ошибка}
Последний верный шаг --- представление правой части как $2^4$. Первый неполный шаг --- запись единственного значения $x=8$ после обращения степени $4/3$.

\err{Потерян отрицательный действительный корень.} Возведение числа в чётную степень уничтожает информацию о его знаке. Поэтому из равенства $y^4=16$ следуют два значения: $y=2$ и $y=-2$. Возведение обеих частей исходного уравнения в степень $3/4$ не учитывает эту развилку при $x<0$.

\keyidea{Сначала выдели кубический корень из $x$, а затем реши уравнение с чётной степенью.} Кубический корень определён и для отрицательных чисел.
\subfield{Исправление от последнего верного шага}
Обозначим $y=\sqrt[3]{x}$. Тогда
\[
  (\sqrt[3]{x})^4=16 \quad\Longleftrightarrow\quad y^4=16.
\]
Значит, $y=\pm2$, а для $x$ получаются два значения.
\subfield{Полное правильное решение}
\begin{align*}
 x^{4/3}&=16,\\
 (\sqrt[3]{x})^4&=16,\\
 y^4&=16,\qquad y=\sqrt[3]{x},\\
 y&=2\quad\text{или}\quad y=-2,\\
 x=y^3&=8\quad\text{или}\quad x=-8.
\end{align*}
\subfield{Проверка}
Для обоих найденных значений имеем
\[
 (\sqrt[3]{8})^4=2^4=16,\qquad
 (\sqrt[3]{-8})^4=(-2)^4=16.
\]
\result{Оба корня подходят: $\boxed{x=-8;\ x=8}$.}
\subfield{Задача на закрепление}
Реши в действительных числах уравнение $x^{4/3}=81$.
\clearpage

\taskhead{12}{Логарифмы: деление на основание}
\subfield{Текст задания}
Реши уравнение
\[
 \log_3(2x-1)=\log_3(x+5)-1.
\]
\subfield{Твоё решение и ответ}
Правильно использовано равенство $1=\log_3 3$ и получено
\[
 \log_3(2x-1)=\log_3\frac{x+5}{3}.
\]
Далее в записи появляется $3(2x-1)=3(x+5)$, откуда найдено $x=6$. \textbf{Твой ответ:} $6$.
\subfield{Где именно возникает ошибка}
Последний правильный шаг --- равенство логарифмов с аргументами $2x-1$ и $(x+5)/3$. Первый неверный шаг --- умножение на $3$ сразу двух выражений, включая числитель правой части.

\err{Неправильно устранён знаменатель.} Из равенства $A=B/3$ при умножении обеих частей на $3$ следует $3A=B$. Получать $3A=3B$ нельзя: знаменатель справа уже сокращается с множителем $3$.

\keyidea{Сначала сравни аргументы логарифмов, затем аккуратно избавься от знаменателя.} Отдельно проверь область допустимых значений.
\subfield{Исправление от последнего верного шага}
\[
 2x-1=\frac{x+5}{3}
 \quad\Longrightarrow\quad
 3(2x-1)=x+5.
\]
\subfield{Полное правильное решение}
Условия существования логарифмов:
\[
 2x-1>0,\qquad x+5>0
 \quad\Longrightarrow\quad x>\frac12.
\]
Поскольку основания логарифмов равны и больше единицы,
\begin{align*}
 \log_3(2x-1)&=\log_3\frac{x+5}{3},\\
 2x-1&=\frac{x+5}{3},\\
 6x-3&=x+5,\\
 5x&=8,\\
 x&=\frac85.
\end{align*}
\subfield{Проверка}
$\frac85>\frac12$. Кроме того,
\[
 2\cdot\frac85-1=\frac{11}{5},\qquad
 \frac{\frac85+5}{3}=\frac{11}{5}.
\]
Аргументы логарифмов совпадают, исходное равенство выполнено.

\result{$\boxed{x=\frac85}$.}
\subfield{Задача на закрепление}
Реши уравнение $\log_2(3x-1)=\log_2(x+7)-1$.
\clearpage

\taskhead{17}{Степенная зависимость температуры}
\subfield{Текст задания}
Мощность излучения определяется формулой $P=\sigma ST^4$, где $\sigma>0$, $S>0$ и температура $T>0$. При неизменных $\sigma$ и $S$ мощность возросла в $\frac{81}{16}$ раза. Во сколько раз изменилась температура?
\subfield{Твоё решение и ответ}
Записана исходная формула $P=\sigma ST^4$, но дальнейшее вычисление и окончательный ответ отсутствуют. \textbf{Твой ответ:} не указан.
\subfield{Где именно решение обрывается}
Последний правильный шаг --- запись зависимости мощности от четвёртой степени температуры. Недостающий шаг --- сравнение двух состояний и сокращение неизменных множителей $\sigma S$.

\err{Решение не завершено.} По одной формуле без отношения конечного и начального значений нельзя получить требуемый коэффициент изменения температуры.

\keyidea{Составь отношение конечной мощности к начальной.} При неизменных коэффициентах оно равно четвёртой степени отношения температур.
\subfield{Исправление от последнего верного шага}
Пусть $P_1,T_1$ --- первоначальные значения, $P_2,T_2$ --- новые. Запишем
\[
 P_1=\sigma ST_1^4,\qquad P_2=\sigma ST_2^4.
\]
Разделим второе равенство на первое, используя $\sigma S>0$ и $T_1>0$.
\subfield{Полное правильное решение}
\begin{align*}
 \frac{P_2}{P_1}&=\frac{\sigma ST_2^4}{\sigma ST_1^4}
 =\left(\frac{T_2}{T_1}\right)^4,\\
 \left(\frac{T_2}{T_1}\right)^4&=\frac{81}{16}
 =\frac{3^4}{2^4}=\left(\frac32\right)^4.
\end{align*}
Температуры положительны, поэтому $\frac{T_2}{T_1}>0$. Из равенства четвёртых степеней следует
\[
 \frac{T_2}{T_1}=\frac32.
\]
\subfield{Проверка}
При увеличении температуры в $\frac32$ раза мощность увеличивается в
$\left(\frac32\right)^4=\frac{81}{16}$ раза, как указано в условии.

\result{Температура увеличилась в $\boxed{\frac32}$ раза.}
\subfield{Задача на закрепление}
При тех же постоянных коэффициентах мощность выросла в $\frac{625}{256}$ раза. Во сколько раз выросла положительная температура?
\clearpage

\taskhead{18}{Логарифмы и квадратное уравнение}
\subfield{Текст задания}
Реши уравнение
\[
 \log_2(x-2)+\log_2(x-5)=3.
\]
\subfield{Твоё решение и ответ}
Верно найдена область допустимых значений $x>5$; затем получены
\[
 (x-2)(x-5)=8,\qquad x^2-7x+2=0,\qquad D=41.
\]
В формуле корней вместо числа $7$ записано $4$. В итоговом ответе указано
$\frac{4+\sqrt{41}}2$.
\subfield{Где именно возникает ошибка}
Последний правильный шаг --- $x^2-7x+2=0$ и $D=41$. Первый неверный шаг --- подстановка коэффициентов в формулу корней квадратного уравнения.

\err{Неверно взят коэффициент при $x$.} Для $ax^2+bx+c=0$ используется формула
\[
 x=\frac{-b\pm\sqrt{b^2-4ac}}{2a}.
\]
В нашем уравнении $a=1$, $b=-7$, $c=2$, следовательно, $-b=7$. Число $4$ в числителе не возникает.

\keyidea{После вычисления дискриминанта ещё раз выпиши $a$, $b$, $c$, а затем подставляй их в формулу.} Найденные корни необходимо проверить по ОДЗ.
\subfield{Исправление от последнего верного шага}
\[
 x^2-7x+2=0
 \quad\Longrightarrow\quad
 x=\frac{7\pm\sqrt{41}}2.
\]
\subfield{Полное правильное решение}
Оба аргумента логарифмов положительны только при $x>5$. Поэтому
\begin{align*}
 \log_2\bigl((x-2)(x-5)\bigr)&=3,\\
 (x-2)(x-5)&=8,\\
 x^2-7x+10&=8,\\
 x^2-7x+2&=0,\\
 D&=(-7)^2-4\cdot1\cdot2=41,\\
 x_{1,2}&=\frac{7\pm\sqrt{41}}2.
\end{align*}
Так как $\sqrt{41}>3$, число $\frac{7-\sqrt{41}}2<2$, значит, оно не удовлетворяет условию $x>5$. А поскольку $\sqrt{41}>3$, число $\frac{7+\sqrt{41}}2>5$ подходит по ОДЗ.
\subfield{Проверка}
Допустимый корень удовлетворяет уравнению $(x-2)(x-5)=8$, поэтому произведение аргументов логарифмов равно $8$, а сумма логарифмов равна $\log_2 8=3$.

\result{Единственный корень: $\boxed{x=\frac{7+\sqrt{41}}2}$.}
\subfield{Задача на закрепление}
Реши уравнение $\log_2(x-3)+\log_2(x-6)=4$.
\clearpage

\taskhead{20}{Знак синуса и двойной угол}
\subfield{Текст задания}
Известно, что
\[
 \cos t=\frac35,\qquad \frac{3\pi}{2}<t<2\pi.
\]
Найди $\sin2t$ и $\cos2t$.
\subfield{Твоё решение и ответ}
Правильно найдено $\sin^2t=\frac{16}{25}$, но затем выбрано $\sin t=\frac45$. Поэтому получено $\sin2t=\frac{24}{25}$. Значение $\cos2t=-\frac7{25}$ найдено верно.
\subfield{Где именно возникает ошибка}
Последний правильный шаг:
\[
 \sin^2t=1-\cos^2t=1-\frac9{25}=\frac{16}{25}.
\]
Первый неверный шаг: $\sin t=\frac45$ без учёта положения угла.

\err{При извлечении корня потерян отрицательный знак.} Из $\sin^2t=16/25$ следует $\sin t=\pm4/5$. На интервале $3\pi/2<t<2\pi$ угол лежит в IV четверти, где синус отрицателен, а косинус положителен.

\keyidea{Знак тригонометрической функции определяется четвертью.} Квадрат позволяет найти модуль, знак уточняется отдельно.
\subfield{Исправление от последнего верного шага}
\[
 \sin t=-\sqrt{\frac{16}{25}}=-\frac45.
\]
Подставим это число в формулу синуса двойного угла.
\subfield{Полное правильное решение}
\begin{align*}
 \sin^2t+\cos^2t&=1,\\
 \sin^2t&=1-\left(\frac35\right)^2=\frac{16}{25},\\
 \sin t&=-\frac45\quad\text{(IV четверть)},\\
 \sin2t&=2\sin t\cos t
 =2\cdot\left(-\frac45\right)\cdot\frac35=-\frac{24}{25},\\
 \cos2t&=\cos^2t-\sin^2t
 =\frac9{25}-\frac{16}{25}=-\frac7{25}.
\end{align*}
\subfield{Проверка}
Оба найденных числа согласованы с тождеством:
\[
 \left(-\frac{24}{25}\right)^2+
 \left(-\frac7{25}\right)^2
 =\frac{576+49}{625}=1.
\]
Кроме того, при $t$ в IV четверти угол $2t$ после вычитания $2\pi$ находится в III или IV четверти в зависимости от $t$; рассчитанные знаки соответствуют конкретному значению $\cos t=3/5$.

\result{$\boxed{\sin2t=-\frac{24}{25}},\quad\boxed{\cos2t=-\frac7{25}}$.}
\subfield{Задача на закрепление}
Известно, что $\sin u=-\frac5{13}$ и $\frac{3\pi}{2}<u<2\pi$. Найди $\sin2u$ и $\cos2u$.
\end{document}
